\documentclass[runningheads]{llncs}

% =====================================================================
%  ACCV 2026 / CV4DC workshop submission (Springer LNCS style).
%  Ported from the NeurIPS Creative AI submission; the body text,
%  figures and appendix are unchanged.
%
%  Build:  latexmk -pdf autopolish.tex        (paper + appendix)
%          latexmk -pdf autopolish_main.tex   (main paper only)
% =====================================================================

% TODO REVIEW: replace '*****' with the OpenReview submission number.
% The [review] option anonymises the author block, adds the paper ID to
% every page, and turns on line numbers.
\usepackage[review,year=2026,ID=*****]{accv}
% TODO FINAL: for the camera-ready version use the line below instead.
%\usepackage{accv}

% Commonly used abbreviations (\eg, \ie, \etc, \cf, \etal).
\usepackage{accvabbrv}

% ---------------------------------------------------------------
% Other packages.  accv.sty already loads: fontenc(T1), amsmath, amssymb,
% xcolor, url, cite, caption, subcaption, lineno -- do not load those again.
\usepackage[utf8]{inputenc}
% Optional (recommended by ACCV): improves PDF readability for those with
% disabilities.  Left off because on this TeX installation it forces PDF 1.6
% and inflates the file from ~6 MB to ~27 MB, over ACCV's 10 MB limit.
% If enabled it must stay above tikz/tcolorbox, which write PDF objects at load.
%\usepackage[accsupp]{axessibility}
\usepackage{graphicx}
\usepackage{booktabs}       % professional-quality tables
\usepackage{tabularx}       % width-constrained wrapping tables
\usepackage{multirow}       % header cells spanning two rows
\usepackage{amsfonts}       % blackboard math symbols
\usepackage{nicefrac}       % compact symbols for 1/2, etc.
\usepackage{microtype}      % microtypography
\usepackage{enumitem}       % compact contribution bullets
\usepackage{tikz}
\usetikzlibrary{arrows.meta, positioning, shapes.geometric, calc, backgrounds, fit}
\usepackage[most]{tcolorbox}

% ---------------------------------------------------------------
% Hyperref (must come after the packages above; accv.sty loads cleveref
% at the end of the preamble, i.e. after hyperref).
% TODO FINAL: comment out the next line and un-comment the one after it.
\usepackage[pagebackref,breaklinks,colorlinks,citecolor=accvblue]{hyperref}
%\usepackage{hyperref}

% TODO FINAL: ORCID icons for the camera-ready author block.
\usepackage{orcidlink}

% ---------------------------------------------------------------
% natbib is not available under LNCS; the LNCS numeric \cite is used
% instead.  Kept as an alias so any stray \citep still compiles.
\providecommand{\citep}{\cite}

% Boxes for prompts / persona reactions / edit instructions (appendix).
\newcommand{\boxfont}{\fontsize{7.5}{9}\selectfont}  % llncs aliases \footnotesize to \small
\tcbset{apbox/.style={boxrule=0.4pt, arc=1.5pt, boxsep=1pt, left=3pt, right=3pt,
  top=1.5pt, bottom=1.5pt, toptitle=0.5pt, bottomtitle=0.5pt, before skip=4pt, after skip=4pt,
  fontupper=\boxfont, fonttitle=\boxfont\bfseries}}
\newtcolorbox{promptbox}[1]{apbox, colback=black!5, colframe=black!80, coltitle=white,
  colbacktitle=black!80, title={#1}}
\newtcolorbox{personabox}[1]{apbox, colback=black!5, colframe=black!80, coltitle=white,
  colbacktitle=black!80, title={#1}, before skip=2pt, after skip=2pt}
\newtcolorbox{instrbox}[1]{apbox, colback=black!5, colframe=black!80, coltitle=white,
  colbacktitle=black!80, title={#1}}
% typewriter prompt text: straight quotes and real hyphens (no ligatures)
\newcommand{\pt}[1]{{\ttfamily\footnotesize #1}}
% Appendix references: a live \ref in the full build; the main-only build has no
% appendix, so it prints the given section number instead of "??".
\newcommand{\appref}[2]{\ifdefined\mainonly #2\else\ref{#1}\fi}

% Palette shared with all figures (Okabe-Ito navy / magenta / yellow).
\definecolor{aenavy}{HTML}{0072B2}     % primary   — synthetic audience
\definecolor{aemagenta}{HTML}{CC79A7}  % secondary — editing loop
\definecolor{aeyellow}{HTML}{F0E442}   % tertiary  — current image (fills only)
\definecolor{aeyellowdk}{HTML}{8C8200} % yellow's strokes: plain yellow vanishes on white
\definecolor{aegray}{HTML}{8A8A86}

% Model-provider marks used as corner badges in Figure 1 (about 7pt across).
% Qwen mark
\newcommand{\qwenmark}{\raisebox{-0.75ex}{\includegraphics[height=7pt]{figs/qwen_logo.png}}}
% Black Forest Labs mark (FLUX), cropped to the glyph; wider than tall, so it is
% set a touch shorter than \qwenmark to read as the same size
\newcommand{\fluxmark}{\raisebox{-0.55ex}{\includegraphics[height=6pt]{figs/flux_logo.png}}}
% held-out judge: DINOv2 (Meta) identity check
\newcommand{\dinomark}{\raisebox{-0.45ex}{\includegraphics[height=5pt]{figs/dinov2_logo.png}}}

% Float packing. The defaults reserve 20%% of every page for text and refuse a
% page that is more than 70%% float, which with this paper's figure load forces
% whole float-only pages and leaves gaps. These let a page be mostly float as
% long as it still carries text, so figures sit beside the prose instead.
\renewcommand{\topfraction}{0.92}
\renewcommand{\bottomfraction}{0.70}
\renewcommand{\textfraction}{0.07}
\renewcommand{\floatpagefraction}{0.88}
\setcounter{topnumber}{3}
\setcounter{bottomnumber}{2}
\setcounter{totalnumber}{5}

\begin{document}

\title{AutoPolish: Improving Visual Aesthetics with Synthetic Audience Critique}

% Shown in the running head of the camera-ready version.
\titlerunning{AutoPolish: Synthetic Audience Critique}

% TODO FINAL: add ORCIDs.  Under the [review] option accv.sty replaces the
% three blocks below with the anonymous placeholder, so they are safe to keep.
\author{Amin Fadaeinejad\inst{1}\thanks{Equal contribution.} \and
Sajjad Pakdamansavoji\inst{1}$^{\star}$\thanks{Corresponding author.}}

\authorrunning{A.~Fadaeinejad and S.~Pakdamansavoji}

\institute{Independent Research\\
\email{aminfadaeinejad.edu@gmail.com}, \email{sj.pakdaman.edu@gmail.com}}

\maketitle

\begin{abstract}
Editing an image with today's text-guided tools is a back-and-forth in which a person decides what should change and an AI system makes the change. Because a person must take part in every round, the process does not scale, and the result follows that one person's taste, which may not suit the people the image is meant for. We propose \textbf{AutoPolish}, which replaces that single opinion with a synthetic audience, many copies of a vision-language model that each react as a different viewer described by a short profile. Their combined feedback stands in for how the intended audience would respond, and an auto-research style loop turns it into edits over a fixed number of rounds, keeping each edit only when an independent scorer rates it higher. Across three datasets of photographs and paintings, we show that only 19 to 47\% of one person's rating is shared with other viewers, while the averages of two separate groups agree at a correlation of up to $0.96$, so one viewer is a poor guide to a crowd. We also show that our audience captures how viewer groups differ up to $10\times$ better than the same model without profiles, and picks the crowd's favorite of two AI-generated images 6 percentage points more often than a simple baseline. In editing, AutoPolish improves images more than a single generic critic on average while keeping them recognizable, and proposes $3\times$ as many distinct changes.

\keywords{Image editing \and Auto-Research \and Persona Simulation \and VLM}
\end{abstract}


\section{Introduction}

For most of its history, image editing meant an operator working through software that offered a fixed set of hand-crafted operations, such as cropping, color curves and masks, so the craft lay largely in mastering the tool. Editors that follow instructions written in plain language have changed this \cite{brooks2023instructpix2pix, flux2024}, because anyone can now ask for an edit in words, and what separates one result from another is no longer skill with the software but the taste of the person giving the instructions. A typical pipeline today is therefore a back-and-forth between an operator and the model, in which the operator describes a change, looks at the result, and keeps asking for further changes until the image looks right to them. This way of working does not scale, because each operator can spend only so much time in the cycle, and it does not serve its audience well, because the final image reflects one person's taste rather than the taste of the people it is meant for.

Two lines of work have tried to take the operator out of this cycle, and each falls short on one of these limits. Computational aesthetics replaces the operator's judgment with a model that scores how good an image looks \cite{murray2012ava, talebi2018nima, ke2023vila, wu2024qalign}, but nearly every such model returns a single score per image, which averages over viewers who often disagree. That disagreement reflects real differences in taste rather than noise \cite{uma2021, plank2022, yang2022para}, so a single score speaks for an average viewer who does not exist, and models of individual taste \cite{ren2017personalized, zhu2022meta, pamela2026} need each user's own ratings. Aesthetic quality is therefore better understood as the reaction of an audience than as a property of the image, and a critic that stands in for the operator has to predict how a group of viewers would respond. Automated pipelines remove the time limit, but they either raise a fixed score in a single training run \cite{xu2023imagereward, black2024ddpo, wallace2024dpo} or let one model critique and repair its own output \cite{madaan2023selfrefine, wu2024sld, wang2024genartist}, in which case the model that proposes an edit also judges it and can mistake its own bias for improvement \cite{xu2024selfbias}. Editing, however, is a repeated cycle of editing, judging and trying again, and an automated editor has to run that cycle itself, trying different edits and keeping only those that a separate judge finds better while the image still looks like the original.

We build \textbf{AutoPolish} around these two ideas. To replace the single score, we create a synthetic audience by asking one frozen vision-language model to react to an image many times, each time as a different viewer described by a short text profile that we call a persona, and we combine these reactions into a collective verdict and a list of complaints. To replace the operator's back-and-forth, we place this audience inside an auto-research style loop~\cite{karpathy2026autoresearch}, that is, a loop in which an agent proposes a change, the change is scored on a fixed metric, and it is kept only if the score improves. In our loop, the audience's complaints become an editing instruction, an image editor carries it out, and an independent scorer from a different model family keeps the edit only if it rates the result higher and the image still looks like the original. Our evaluations show that the audience suggests three times as many distinct edits as a single generic critic, such as adding plants and artwork to a bare office, while also improving the image more on average. A single simulated viewer barely beats a coin flip at guessing which of two images people prefer, yet twenty of them together agree with people two times in three, which is why we rely on an audience rather than one voice.

Our contributions are as follows.

\begin{itemize}[leftmargin=1.4em, itemsep=2pt, topsep=2pt, parsep=0pt]
\item We simulate collective taste with a synthetic audience that captures how viewer groups differ $10\times$ better than the same model without profiles.

\item We model the back-and-forth of editing with a custom auto-research loop that beats a generic critic on average and suggests $3\times$ as many distinct edits.

\item We evaluate on photographs, paintings and AI-generated images, where the audience beats a simple baseline by 6 percentage points on AI-generated pairs.
\end{itemize}

\begin{figure}[t]
  \centering
    \begin{tikzpicture}[
    scale=0.82, transform shape,
    font=\small,
    box/.style={rectangle, rounded corners=3pt, draw=aegray!60, thick, align=center,
                inner sep=3pt, minimum height=11mm, minimum width=27mm, fill=white},
    aud/.style={box, draw=aenavy, fill=aenavy!7},
    edt/.style={box, draw=aemagenta, fill=aemagenta!10},
    jdg/.style={box, draw=aemagenta, fill=aemagenta!10},
    io/.style={box, draw=aegray!70, fill=aegray!8},
    img/.style={box, draw=aeyellowdk, fill=aeyellow!30},
    flow/.style={-{Latex[length=2.4mm]}, thick, draw=aegray!85},
    feed/.style={-{Latex[length=2.4mm]}, thick, draw=aenavy, dashed},
    look/.style={-{Latex[length=2.4mm]}, thick, draw=aeyellowdk, dashed},
    ret/.style={-{Latex[length=2.4mm]}, thick, draw=aemagenta!80!black},
    note/.style={font=\scriptsize\itshape, text=aegray},
  ]
    % ===== Top band: the synthetic audience (4 aligned columns) =====
    % persona-card stack: two offset cards behind the labelled one, so the
    % block reads as N cards rather than one
    \node[io, draw=aegray!40] (pc2) at (1.16,0.22) {};
    \node[io, draw=aegray!55] (pc1) at (1.03,0.10) {};
    \node[io] (pers) at (0.9,-0.03)
      {Persona cards\\[-2pt]{\scriptsize (age, nationality, etc.)}\\[-4pt]{\scriptsize\strut}};
    \node[anchor=south east, inner sep=2pt] at (pers.south east) {\scriptsize $\times N$};
    \node[aud] (vlm)  at (4.5,-0.03) {Role-play\\[-2pt]{\scriptsize Qwen2-VL-7B}};
    \node[anchor=south east, inner sep=2pt] at (vlm.south east) {\qwenmark};
    % what the role-play emits: one reaction per card, so it mirrors the input
    % stack -- gray artifact, stacked, times N
    \node[io, draw=aegray!40] (rc2) at (8.36,0.22) {};
    \node[io, draw=aegray!55] (rc1) at (8.23,0.10) {};
    \node[io] (resp) at (8.1,-0.03)
      {Reactions\\[-2pt]{\scriptsize (score, complaint, etc.)}\\[-4pt]{\scriptsize\strut}};
    \node[anchor=south east, inner sep=2pt] at (resp.south east) {\scriptsize $\times N$};
    \node[aud] (agg)  at (11.7,-0.03){Aggregate\\[-1pt]group reaction\\[-2pt]{\scriptsize (stats, complaints, etc.)}};

    \draw[flow] (pers) -- (vlm);
    \draw[flow] (vlm)  -- (resp);
    \draw[flow] (resp) -- (agg);

    \node[draw=aenavy!55, dashed, rounded corners=5pt, inner sep=7pt,
          fit=(pers)(pc1)(pc2)(vlm)(resp)(rc1)(rc2)(agg)] (frA) {};
    \node[anchor=south west, text=aenavy, font=\footnotesize\bfseries]
      at (frA.north west) {Synthetic audience};

    % ===== Bottom band: the AutoPolish loop (same 4 columns) =====
    \node[img] (cur)  at (0.9,-3.05)  {Current\\best image};
    \node[edt] (crit) at (4.5,-3.05)  {Distill instruction\\[-2pt]{\scriptsize Qwen2-VL-7B}};
    \node[anchor=south east, inner sep=2pt] at (crit.south east) {\qwenmark};
    \node[edt] (ed)   at (8.1,-3.05)  {Image editor\\[-2pt]{\scriptsize FLUX.1-Kontext}};
    \node[anchor=south east, inner sep=2pt] at (ed.south east) {\fluxmark};
    \node[jdg] (jud)  at (11.7,-3.05) {Held-out judge\\[-2pt]{\scriptsize LAION $+$ DINOv2}};
    \node[anchor=south east, inner sep=2pt] at (jud.south east) {\dinomark};

    \draw[flow] (cur)  -- (crit);
    \draw[flow] (crit) -- (ed);
    \draw[flow] (ed)   -- (jud);

    \node[draw=aemagenta!60, dashed, rounded corners=5pt, inner sep=7pt,
          fit=(cur)(crit)(ed)(jud)] (frB) {};
    \node[anchor=south west, text=aemagenta!80!black, font=\footnotesize\bfseries]
      at (frB.north west) {Auto-research loop};

    % the panel looks at the image the loop is currently holding. Routed as a
    % Manhattan path, not a diagonal: a diagonal here crosses its own label and
    % the feedback arrow's lane.
    \draw[look] (2.7,-3.05) -- (2.7,-1.35) -| (vlm.south);
    \fill[aegray!85] (2.7,-3.05) circle (1.5pt);

    % feedback: the audience's complaints steer the critic
    \draw[feed] (agg.south) |- (4.5,-1.80) -- (crit.north);
    \node[note, text=aenavy, anchor=east] at (10.5,-1.62) {complaints steer the instruction};

    % accept-if-better return: judge commits back to the current best
    \draw[ret] (jud.east) -- ++(0.7,0) -- ++(0,-1.40) -| (cur.south);
    % one line in the gap between the dashed loop frame and the return arrow
    \node[note, text=aemagenta!70!black] at (6.3,-4.14)
      {keep the edit only if it scores higher \ $\cdot$ \ $10$ steps, always from the original};
  \end{tikzpicture}
  \captionsetup{skip=3pt}% tighter image-to-caption gap, this float only
  \caption{\textbf{Overview of AutoPolish.} \textbf{(Top)} A frozen VLM role-plays $N$ plain-text persona cards, and their reactions to an image are aggregated into a group critique. \textbf{(Bottom)} The VLM distills this critique into one edit instruction for FLUX.1-Kontext, and a held-out judge keeps the edit only if its aesthetic score rises while DINOv2 similarity to the original stays high; otherwise it rolls back.}
  \label{fig:overview}
\end{figure}


\section{Related work}


\subsubsection{Vision-language models as judges.} Using another AI model as a judge has become a widely adopted approach for evaluating the output of a model \cite{zheng2023judging, liu2023geval}. For aesthetics, the early zero-shot scorers were built on CLIP embeddings \cite{wang2023clipiqa}, which made open-vocabulary image quality assessment practical, and were followed by aesthetic predictors trained on CLIP embeddings \cite{schuhmann2022} or pretrained on image-comment pairs \cite{ke2023vila}. Later, multimodal LLMs were used as scorers by deriving a continuous score from the logits of the quality tokens \cite{wu2024qbench}, and instruction tuning them to predict discrete text-defined quality levels instead of a numerical score improved the alignment with human ratings \cite{wu2024qalign}. However, general-purpose VLMs are still weaker than humans in detailed aesthetic reasoning and explanation \cite{huang2024aesbench, zhou2024uniaa}, which motivated models fine-tuned specifically for aesthetic perception \cite{zhou2024uniaa, huang2024aesexpert}. None of these approaches is conditioned on the viewer: each of them models one consensus judgment per image, typically a single general Mean Opinion Score (MOS). As a result, the disagreement inside the audience is collapsed into a static score, and there is no mechanism to model how aesthetic judgments vary across viewer groups or demographics.

\subsubsection{Simulating human behavior with language models.} Language models conditioned on a person's demographics have been shown to reproduce human survey answers with surprising fidelity \cite{argyle2023}, suggesting that persona-conditioned models can serve as a proxy for human samples. Later work, however, reports clear limitations: simulated answers vary less than real ones \cite{bisbee2024synthetic}, several groups are misrepresented \cite{santurkar2023whose}, and persona prompts slide easily into stereotype \cite{cheng2023compost}. Together, these findings indicate that such simulations are trustworthy at the level of groups rather than individuals, and only once calibrated. Moreover, this literature is almost entirely about text, which leaves open whether a persona also changes how a model perceives an image. We address it for aesthetic judgment.

\subsubsection{Bias in aesthetic models.} Aesthetic predictors are used to filter web-scale training corpora \cite{schuhmann2022}, so their bias passes to the models trained on the filtered data. Taylor et al.\ \cite{taylor2026gaze} audit the LAION aesthetic predictor and find that it gives the highest scores to realistic landscapes, cityscapes and portraits by western and Japanese artists. For multimodal LLMs, AesBiasBench \cite{aesbiasbench2025} measures how aesthetic evaluations vary across demographic groups, and reports that adding identity information to the prompt often increases stereotype bias. 

\subsubsection{Feedback loops for editing.} A line of work improves the outputs of generative models using feedback. Most of these methods rely on a reward, such as a learned preference model \cite{xu2023imagereward, kirstain2023pickapic}, reinforcement learning \cite{black2024ddpo, fan2023dpok}, reward gradients \cite{clark2024draft, prabhudesai2023alignprop}, or preference optimization \cite{wallace2024dpo}. Other methods replace the reward with language feedback, where a model criticizes its own output \cite{madaan2023selfrefine, shinn2023reflexion, gou2024critic}. For images, a vision-language model identifies the problem and a correction is applied to the image \cite{wu2024sld, wang2024genartist}, for example by an instruction-based editor \cite{brooks2023instructpix2pix, flux2024}. These approaches have two limitations. First, a reward is a single number, so the disagreement between viewers is already lost before the model receives it. Second, when the same model proposes an edit and also scores it, the gain it reports can come from its own bias instead of a real improvement \cite{xu2024selfbias}. AutoPolish addresses both. Its feedback comes from a group of simulated viewers and is given as text instead of a score, and the model that judges the result is from a different model family than the one that proposes the edit \cite{schuhmann2022, oquab2023}. The loop itself follows the auto-research pattern~\cite{karpathy2026autoresearch}, in which an agent edits a training script, runs a short experiment and keeps the change only if a validation metric improves, and we apply this propose, score and keep-if-better structure to image editing.


\section{Method}
\label{sec:method}

\subsubsection{Persona cards.} A persona card is a short textual description of a real annotator, generated deterministically from the attributes their dataset provides: age band, education, art and photography familiarity, Big-Five scores, and nationality. Attributes that a dataset does not provide are left out, not filled in with estimated values. A typical card reads \emph{``Viewer profile: 35 to 44, master's degree, high art familiarity, moderate photography experience, nationality: Belgium. Personality (Big-5, 1 to 5): O 4.2, C 3.1, E 2.8, A 3.9, N 2.5.''} We use the same template for the annotators we use from the Synthetic audience. Our claims are therefore limited to the attributes that the card actually contains.
\subsubsection{The frozen LLM judge.} The judge is a pretrained vision-language model that we query but never fine-tune. Each query supplies an image, a persona card, and an instruction to answer in a fixed JSON format; the model returns a score, a short rationale, and a few auxiliary fields. We parse it as strict JSON. If parsing fails, a per-field regular expression recovers the values, so we keep the sample instead of dropping it.  The backbone is \texttt{Qwen2-VL-7B-Instruct} \cite{wang2024qwen2vl}, in bf16. The “persona-blind” control format behaves as if no prompt conditioning were involved for the VLM judgment.
\subsubsection{Score scales.} Each dataset collected its ratings on its own scale, and we ask the judge for a score on that same scale: the prompt states the range and what its ends and middle mean, and the returned value is clamped to the range and rounded to the nearest step the dataset allows. On PARA \cite{yang2022para}, each rater gives an overall aesthetic score from 1 to 5 with a middle option between every two whole numbers, so the score moves in steps of 0.5, and also rates eight further aspects of the photograph, such as quality, composition, color, lighting and how much they like its content, on whole numbers from 1 to 5. On EVA \cite{kang2020eva}, the overall score is a whole number from 0 (least beautiful) to 10 (most beautiful), and light and color, composition and depth, quality, semantics and how hard the image was to judge are each rated on four levels, from very bad to very good. On LAPIS \cite{maerten2025lapis}, raters judge the aesthetic value of a painting on a 0 to 100 slider with seven labelled tick marks, the finest of the three scales and the only rating the dataset collects. Rapidata \cite{rapidata2024} has no scale: each voter sees two images and answers ``Which image do you prefer?'', so a vote is a binary choice, and the score of an image in a pair is the share of voters who prefer it. We study the overall aesthetic score throughout, and the other axes enter only the breadth check of Appendix~\appref{sup:robust}{A5}, which covers all 16 rated axes. Because the scales differ, we rescale every score linearly to $[0,1]$ using the ends of its scale before comparing numbers across datasets, so the errors we report are on this common scale.
\subsubsection{Held-out aesthetic score.} Human ratings exist only for the images in the datasets, not for the edits that AutoPolish produces, so to score an edit we use the LAION aesthetic predictor \cite{schuhmann2022, schuhmann2022predictor}, the model used to select the LAION-Aesthetics subsets of LAION-5B. It embeds an image with a frozen CLIP ViT-L/14 encoder \cite{radford2021}, normalizes the embedding to unit length, and maps it to a single number with a small MLP head. The head was fit by regression on human ratings from three sources, namely the 1 to 10 ratings of AI-generated images in Simulacra Aesthetic Captions \cite{pressman2022sac}, the photograph ratings of AVA \cite{murray2012ava}, and a set of rated logos (LAION-Logos), so its output estimates the average rating that the raters in these sources would give, on a scale of roughly 1 to 10. Unlike the persona-conditioned judge, it is not conditioned on any viewer and reads only the image, so the wording of an edit instruction cannot sway it. The critics in AutoPolish also return a score of their own, from 0 to 100, but we only record it and never use it to accept an edit.
\subsubsection{From a panel to a group prediction.} To predict how a viewer group responds to an image, we sample $N$ persona cards from the dataset and use them to judge the image, and summarize the reactions: the mean and spread of the scores, the most frequent complaints extracted from the rationales, and compare them with the ground truth scoring.

Averaging over the panel is not enough on its own, because the judge's error is systematic and not random. Its ratings are pulled toward the middle of the scale: their standard deviation is only 51-58\% of the human standard deviation, 34-56\% of them land on a single value, and the endpoints are almost never used. Averaging cancels the random part of the error, but the compression is shared across the group population, so it survives. This is why raw population averages are less accurate than simply predicting the population mean. We therefore correct the scores after inference, using isotonic regression \cite{niculescu2005}. This fits a monotone mapping from the judge's raw score to a calibrated one. We fit the mapping with two-fold cross-fitting over images, so every rating is calibrated by a mapping that was estimated without it. Because the mapping is monotone, it changes the absolute error but leaves every ranking as it was, and because it acts only on the output, the model weights stay frozen. The mapping also transfers across datasets. We can therefore fit it once on real images, where human ratings are available, and apply it unchanged to generated images, for which no per-image ratings exist.

\subsubsection{AutoPolish.} AutoPolish uses the group as the critic in an iterative editing loop, under the constraint that the model proposing edits and the model judging them are different, since a loop that scores its own edits can improve its objective without improving the image. Each image is refined for $R=10$ steps of (i) \textbf{critique}: the group views the current image, and its complaints are aggregated across personas and accumulated across steps into a single instruction of at most fifteen words; (ii) \textbf{anchored re-edit}: FLUX.1-Kontext-dev \cite{flux2024} applies that instruction to the \emph{original} image rather than the previous edit, so editing artifacts do not compound, producing $K=3$ candidates; (iii) \textbf{held-out scoring}: each candidate is scored by the LAION aesthetic predictor described above, a different model family from the critic, and checked for identity drift with DINOv2 \cite{oquab2023}; (iv) \textbf{accept-if-better}: the new image replaces the previous best image if the aesthetic score improves and DINOv2 similarity to the original stays above $0.78$, so the best-so-far score increases monotonically. We compare four critics: \textsc{static}, a fixed ``improve this image'' instruction; \textsc{blind}, a single persona-free critique; \textsc{society}, ten aggregated personas, which is our method; and \textsc{reward-only}, which optimizes the aesthetic predictor directly and is an upper reference rather than a comparable baseline. Every condition appends the same suffix asking for a visible edit, so the request sent to the editor is worded the same in every condition, and we deliberately do not score instruction adherence, which would favour whichever condition writes the instruction that is easiest to follow.

\section{Experimental setup}

\subsubsection{Data.} We use four publicly released datasets in which every record ties an aesthetic judgment to the person who made it, together with personal information about that person. PARA \cite{yang2022para} contains photographs, of which we use 4{,}000 with about 52k individual ratings, and it describes its raters in the most detail, recording their age, gender, education, art and photography experience, and their Big-Five personality traits of openness, conscientiousness, extraversion, agreeableness and neuroticism. EVA \cite{kang2020eva} also contains photographs, drawn from AVA \cite{murray2012ava}, with 4{,}070 images and 137k ratings, but it records less about its raters, namely their age, gender, region and photographic skill. LAPIS \cite{maerten2025lapis} contains paintings, of which we use 4{,}000 with 90k ratings, and it records each rater's nationality, interest in art and age. Rapidata \cite{rapidata2024} contains 898 pairs of AI-generated images with 23{,}445 votes, in which each voter picks the image they prefer, and it records the country and language of every voter across 120 countries.

\subsubsection{Models and baselines.} We compare three variants throughout the paper. The blind variant uses the model with no viewer profile, which in editing becomes a single critic or a fixed generic instruction, and it shows what the model does on its own. The synthetic audience, our approach, asks the same model to react to an image many times, each time as a different viewer described by a short profile, and combines these reactions into a group verdict that, in editing, also serves as the critic that guides each round. The oracle has access to the real ratings, either as a bar the synthetic audience must clear, such as always predicting the average rating or the more popular image, or as a ceiling, such as a second group of real viewers or, in editing, a loop that optimizes the scorer directly. The exact form of each variant depends on the experiment, and we describe it alongside the results.

\subsubsection{Evaluation protocol.} To test whether the audience tells viewer groups apart, we use between-group separation. We split the raters of each image into groups by one attribute, such as age or interest in art, take the gap between every pair of groups in both the real and the predicted ratings, and report the correlation between the two, which rewards getting differences between viewers right rather than the overall average. To test whether a model knows which of two images people prefer, we report how often it picks the right one. To measure what editing achieves, we report the gain in aesthetic score and the share of images that improve, while checking that each edited image stays close to the original. We use bootstrap resampling over images to give a 95\% confidence interval for each difference between variants, and we correct for multiple testing when running several group tests. Throughout, the judge's scores are mapped onto the human rating scale by a correction fit on one half of the images and applied to the other, so no score is corrected using its own rating.


\section{Results}

\subsection{Taste is better reproduced for groups than for individuals}
\label{sec:ceiling}

Before testing any model, we ask how predictable aesthetic taste is in the first place, by comparing how reliable one person's rating is with how reliable the average of a group is. We treat each rating as the sum of two parts: the image's average rating, which all of its raters share, and the rater's own deviation from that average. The intraclass correlation (ICC)~\cite{shrout1979} then measures how much of the variation in ratings comes from the shared part. For a single rating, ICC(1) is the expected correlation between two people rating the same images. For a group average, ICC($k$) is the expected correlation between the averages of two separate groups of $k$ people, where $k$ is the number of raters per image, 23 to 34 depending on the dataset. Finally, to see how much of personal taste a profile can explain, we regress each rater's deviations on their profile attributes.

\begin{table}[t]
  \caption{Reliability of a single rating and of a group average, the share of each rater's personal taste that their profile explains, and the average number of raters per image across three datasets. A group average is far more reliable than one rating, and the profile explains little of personal taste.}
  \label{tab:ceiling}
  \centering
  \fontsize{7.5}{9}\selectfont  % llncs aliases \footnotesize to \small
  \setlength{\aboverulesep}{0pt}\setlength{\belowrulesep}{0pt}
  \renewcommand{\arraystretch}{1.0}
  \newcolumntype{C}[1]{>{\hsize=#1\hsize\centering\arraybackslash}X}
  \begin{tabularx}{\linewidth}{l|C{0.92}C{0.92}|C{0.96}|C{1.2}}
    \toprule
    \multirow{2}{*}{Dataset} & \multicolumn{2}{c|}{Reliability $\uparrow$\rule{0pt}{2.4ex}\rule[-0.7ex]{0pt}{0pt}} & \multirow{2}{=}[-0.5ex]{\centering Taste explained\newline by profile} & \multirow{2}{=}{\centering \# Raters/image} \\
    \cline{2-3}
    \rule[-0.8ex]{0pt}{0pt} & single rating\rule{0pt}{2.3ex} & group average & & \\
    \midrule
    PARA~\cite{yang2022para}\rule{0pt}{2.3ex}      & 0.470 & \textbf{0.958} & 2.0\% & 26 \\
    EVA~\cite{kang2020eva}        & 0.223 & \textbf{0.906} & 0.9\% & 34 \\
    LAPIS~\cite{maerten2025lapis}\rule[-0.7ex]{0pt}{0pt} & 0.188 & \textbf{0.839} & 4.0\% & 23 \\
    \midrule
    Average\rule{0pt}{2.3ex}\rule[-0.7ex]{0pt}{0pt} & 0.293 & \textbf{0.901} & 2.3\% & 27 \\
    \bottomrule
  \end{tabularx}
\end{table}

Table~\ref{tab:ceiling} shows a large gap between the two reliabilities. Averaged over the three datasets, a single rating reaches a reliability of only 0.29, so most of what one person says about an image is their own, while the average of a group reaches 0.90, which means a second group of viewers would give nearly the same average. The gap also varies with the kind of image, since a single rating is most reliable on the PARA photographs at 0.47, drops to 0.22 on the EVA photographs, and is lowest on the LAPIS paintings at 0.19, where a group average still holds at 0.84. Paintings therefore draw the most personal reactions, although the difference between the two photograph datasets suggests that the type of image is not the only factor. The profile attributes explain only 2.3\% of personal taste on average and at most 4\%. Personalized reward models that learn an embedding for each user from that user's ratings predict individual taste accurately~\cite{pamela2026}, but they need ratings from every viewer, which a persona card does not contain. Together, these results suggest that, when only a profile is available, the group is a more realistic target than the individual, so we evaluate the synthetic audience on groups in the rest of the paper.

\begin{figure}[t]
  \centering
  \captionsetup{skip=3pt}% tighter image-to-caption gap, this float only
  \includegraphics[width=\linewidth]{figs/pf_persona.png}
  \caption{The audience's rating error and what its comments reveal about the viewer. \textbf{(left)} Error of the audience's average rating before and after calibration, against the prior. \textbf{(middle)} How well a classifier recovers a rater's attribute from the comment alone and \textbf{(right)} distinct comments per rater of an image, both with and without a persona. Calibration beats the prior on every dataset, and the persona makes the comments more varied and more telling of the viewer. Exact values are given in Appendix~\appref{sup:fig2}{A4}.}
  \label{fig:persona}

  \vspace{6pt}
  \includegraphics[width=\linewidth]{figs/pf_steer.png}
  \caption{How far each group of viewers departs from the crowd in its real ratings (horizontal) and in the VLM's ratings under that group's persona (vertical), on each dataset. Each point is one group, marker size shows how many ratings it has, and the dashed line marks perfect agreement. The VLM follows the real groups on PARA and LAPIS but not on EVA.}
  \label{fig:steer}
\end{figure}

\subsection{Personas change how the VLM reacts to an image}
\label{sec:persona}

We compare the reactions of a persona-conditioned VLM with those of the blind variant, asking how the two differ and which is closer to how real viewers react. First, we compare the audience's calibrated average for each image with the real average and with a prior that always predicts the overall average rating. Second, to check whether the persona moves the scores in the right direction, we subtract each image's average from every rating, which leaves only how each rater departs from the crowd, then average these departures over each group of viewers, such as people with more or less art experience, and correlate the VLM's group departures with the real ones. Third, we train a logistic regression on TF-IDF features to guess a rater's attribute from their comment alone, using art experience on PARA, photographic skill on EVA and interest in art on LAPIS, and we count how many distinct comments the VLM writes for the same image. Training details and detailed results are provided in Appendix~\appref{sup:probe}{A3}.

Figure~\ref{fig:persona} (left) shows that the raw scores of the synthetic audience are not directly usable. The VLM rates almost every image near the middle of the scale, so on every dataset the audience's average lands further from the real average than the prior does. A single calibration fixes this, cutting the audience's error roughly in half, from 0.14 to 0.07 on average against 0.10 for the prior, and it beats the prior on each dataset separately.

Figure~\ref{fig:steer} shows that the persona moves the VLM's scores in the same direction that real viewers differ. On PARA and LAPIS, which describe their raters in detail, the VLM's group departures follow the real ones with a correlation of 0.37 and 0.39. On EVA they do not, with a negative correlation of $-0.24$, which we attribute to EVA recording little about its raters and using a coarse whole-number scale that leaves little room for small shifts in taste.

The impact of the persona is clearer in the written comments than in the scores. A classifier trained on the comments alone tells viewer groups apart well above chance on all three datasets, scoring 0.58 to 0.70, while on the blind variant it stays at chance (Figure~\ref{fig:persona}, middle). The persona also makes the comments far more varied. On average, the VLM writes about 37 distinct comments per 100 raters of an image, nearly four times the 10 of the blind variant (Figure~\ref{fig:persona}, right). Together, these results suggest that the persona does change how the VLM reacts to an image, a little in the scores it gives and much more in what it writes.

\subsection{The synthetic audience simulates population reactions}
\label{sec:audience}
\label{sec:generated}

We test whether a panel of persona-conditioned VLMs, the synthetic audience, can stand in for the reaction of a target population, at two levels. At the level of distributions, we ask whether panels built for different groups of viewers reproduce the real differences between those groups. Using the between-group separation of the evaluation protocol, we split the raters of each image by one attribute, namely art experience, photography experience and age on PARA, photographic skill and age on EVA, and nationality, interest in art and age on LAPIS, keep groups with at least four raters, and compare the predicted gaps between groups with the real ones, which gives eight tests across the three datasets. At the level of a task, we ask whether the panel picks the same image as the crowd. On the AI-generated pairs of Rapidata, the VLM casts one vote for every human vote, conditioned on that voter's country and language. For each pair we draw a panel of one to twenty of these votes, take its majority choice and compare it with the crowd's majority, averaging over 200 random draws, against a baseline that always picks the image from the model the crowd prefers overall.

\begin{figure}[t]
  \centering
  \begin{minipage}[t]{0.52\linewidth}
    \vspace{0pt}
    \centering
    \captionof{table}{Separation between viewer groups for each attribute, with $p$-values corrected for running eight tests (Holm, $\alpha=0.05$). The audience gets the direction right in five of the eight cases and backwards for age on EVA.}
    \label{tab:holm}
    \vspace{0pt}
    \scriptsize
    \setlength{\tabcolsep}{2.2pt}
    \renewcommand{\arraystretch}{0.82}
    \begin{tabular}{llccc}
      \toprule
      Dataset & slice & sep.\ $r\,\uparrow$ & Holm $p\,\downarrow$ & surv. \\
      \midrule
      LAPIS & age                   & $+0.102$ & $<10^{-3}$ & \checkmark \\
      LAPIS & art interest          & $+0.088$ & $<10^{-3}$ & \checkmark \\
      LAPIS & nationality           & $+0.068$ & $8\times10^{-5}$ & \checkmark \\
      PARA  & photo.\ exp.\     & $+0.055$ & $1.1\times10^{-4}$ & \checkmark \\
      PARA  & art exp.\             & $+0.040$ & $4.2\times10^{-3}$ & \checkmark \\
      PARA  & age                   & $+0.008$ & $0.50$ & $\times$ \\
      EVA   & photo.\ level    & $-0.018$ & $0.50$ & $\times$ \\
      EVA   & age                   & $-0.079$ & $<10^{-3}$ & \checkmark \\
      \bottomrule
    \end{tabular}
  \end{minipage}\hfill
  \begin{minipage}[t]{0.45\linewidth}
    \vspace{0pt}
    \centering
    \includegraphics[width=\linewidth]{figs/pf_separation.png}
    \caption{Separation between viewer groups on each dataset and on average, with 95\% bootstrap intervals. The persona gets the direction right on LAPIS and PARA but not on EVA, while the blind variant stays near zero.}
    \label{fig:audience}
  \end{minipage}
\end{figure}

Figure~\ref{fig:audience} shows how well the predicted differences between groups follow the real ones, where a positive value means the audience gets the direction of the differences right. On LAPIS it does so clearly, with a correlation of 0.166 against 0.016 for the blind variant. On PARA the effect is weaker but still present (0.044), while on EVA the predictions point the opposite way ($-0.084$). Averaged over the three datasets, the audience reaches 0.042 against 0.003 for the blind variant. Table~\ref{tab:holm} breaks these results down by the attribute used to form the groups. In five of the eight cases the audience gets the direction right. On EVA, grouping viewers by age gives the opposite direction, which explains the negative result for EVA in Figure~\ref{fig:audience}, and the remaining two cases show no clear effect either way. The order LAPIS, PARA, EVA follows how much each dataset records about its raters, which suggests that the audience reproduces the differences within a population as far as its profiles describe them.

\begin{figure}[t]
  \centering
  \captionsetup{skip=3pt}% tighter image-to-caption gap, this float only
  \includegraphics[width=\linewidth]{figs/pf_generated.png}
  \caption{The synthetic audience on AI-generated image pairs. \textbf{(left)} Agreement with the crowd's majority as the panel grows, against the prior of always picking the more popular model. \textbf{(right)} Panel against crowd preference per pair, colored by number of voters. Agreement passes the prior with more personas, and the panel tracks the crowd.}
  \label{fig:generated}
\end{figure}

Figure~\ref{fig:generated} (left) shows that the vote of a single persona-conditioned VLM agrees with the crowd's majority on 58.8\% of pairs, just below the 60.0\% of always picking the more popular model, but agreement rises as the panel grows and reaches 65.7\% with twenty personas. Figure~\ref{fig:generated} (right) shows that, per pair, the share of the panel that prefers an image tracks the share of the crowd that does, with a correlation of 0.47, so the panel gives a graded preference rather than a lucky guess. Together, these results suggest that the votes of many persona-conditioned VLMs add up to a usable estimate of a population's taste.

\subsection{The synthetic audience can steer image editing}
\label{sec:c4}

We test whether the synthetic audience can be used for automated image editing as the critic in the loop of AutoPolish, improving images while keeping them recognizable. We run the loop for ten rounds on 100 photographs, 50 from PARA and 50 from EVA, and compare four critics that differ only in where the editing instruction comes from. The first two critics are blind variants, namely a fixed instruction to improve the image (static string) and a single VLM critique with no persona (blind VLM). The third is AutoPolish, which distills the complaints of ten personas into one instruction, and the fourth is the reward-only ceiling, which optimizes the independent scorer directly and serves as an upper reference rather than a fair competitor. All critics end their instruction with the same request for a visible edit (see Appendix~\appref{sup:prompts}{A8}). All scores come from the independent scorer, a LAION aesthetic predictor that rates images on a scale of roughly 1 to 10. Besides the gain in score and the share of images that improve, we report the gain on the images that do improve, the edit size, which is one minus the DINOv2 similarity between the final image and the original and is capped at 0.22 by the similarity floor of 0.78 that the loop enforces, and how many distinct complaints the critic raises per round.

\begin{table}[b]
  \caption{Results of each critic over ten rounds on 100 photographs, where edit size is how much the final image differs from the original. AutoPolish gives the largest gain among the critics, the largest gain on the images it improves, and the most feedback.}
  \label{tab:c4}
  \centering
  \fontsize{8}{9.5}\selectfont
  \setlength{\aboverulesep}{0pt}\setlength{\belowrulesep}{0pt}
  \renewcommand{\arraystretch}{1.1}
  \setlength{\tabcolsep}{3.5pt}
  \begin{tabular*}{\linewidth}{@{\extracolsep{\fill}}l|cccccc@{}}
    \toprule
    \multirow{2}{*}{Critic} & \multirow{2}{*}{Gain $\uparrow$} & \multirow{2}{*}{95\% CI} & \multirow{2}{*}{Improved} & Gain per\rule{0pt}{2.4ex} & Edit & \multirow{2}{*}{Complaints $\uparrow$} \\
    \rule[-0.8ex]{0pt}{0pt} & & & & improved $\uparrow$ & size & \\
    \midrule
    static string\rule{0pt}{2.3ex} & $+0.164$ & [0.135, 0.196] & 84\% & 0.195 & 0.032 & n/a \\
    blind VLM            & $+0.155$ & [0.126, 0.185] & 72\% & 0.215 & 0.036 & 1.00 \\
    \textbf{AutoPolish (ours)}\rule[-0.7ex]{0pt}{0pt} & $\mathbf{+0.170}$ & [0.136, 0.207] & 69\% & \textbf{0.247} & \textbf{0.049} & \textbf{3.12} \\
    \midrule
    \textit{reward-only (ceiling)}\rule{0pt}{2.3ex}\rule[-0.7ex]{0pt}{0pt} & $+0.187$ & [0.157, 0.221] & 86\% & 0.218 & 0.040 & n/a \\
    \bottomrule
  \end{tabular*}
\end{table}

Table~\ref{tab:c4} shows that AutoPolish gives the largest score gain of the three critics, raising the independent score by 0.170 on average, from 5.10 to 5.27, against 0.164 for the static string and 0.155 for the blind VLM. Our method improves fewer images than the static string, 69\% against 84\%, but changes the ones it improves more. The gain on those images is 0.247, higher than for any other critic and the reward-only ceiling (0.218), and its edits are the largest (0.049), while every critic stays well below the cap of 0.22. AutoPolish also raises about three distinct complaints per round, against one for the blind VLM, while the static string always repeats the same instruction.

\begin{figure}[t]
  \centering
  \captionsetup{skip=3pt}% tighter image-to-caption gap, this float only
  \includegraphics[width=\linewidth]{figs/pf_autopolish_accv.png}
  \caption{How each critic behaves over ten rounds. \textbf{(left)} Best score so far at each round, averaged over 100 photographs, with a 95\% bootstrap band. \textbf{(right)} Final gain of each image against its similarity to the original, with the 0.78 floor marked. AutoPolish leads from the fourth round, and almost every image stays far above the floor.}
  \label{fig:c4quant}
\end{figure}

Figure~\ref{fig:c4quant} (left) shows that most of the gain arrives in the first few rounds, but the difference between critics becomes clearer as the rounds go on, with AutoPolish starting slightly behind and pulling ahead of the other critics from the fourth round. Figure~\ref{fig:c4quant} (right) shows that almost every final image stays well above the similarity floor. Figure~\ref{fig:c4prog} shows how AutoPolish gradually improves images over several rounds while keeping them similar to the original. Figure~\ref{fig:c4qual} compares the critics on the same images and shows the larger edits and gains of AutoPolish. Together, these results suggest that the synthetic audience is an effective critic for editing, exceeding simpler critics in average score while making larger and more concrete changes. Appendix~\appref{sup:c4qual}{A7} gives more examples.

\begin{figure}[tb]
  \centering
  \captionsetup{skip=3pt}% tighter image-to-caption gap, this float only
  \includegraphics[width=\linewidth]{figs/pf_progression_accv.png}
  \caption{How AutoPolish improves images over its rounds, shown by the best image so far at selected rounds with its LAION aesthetic score above each image. The right column explains the change from source to last round. Each round adds the complaints gathered so far, so the images improve step by step while staying close to the original.}
  \label{fig:c4prog}
\end{figure}

\begin{figure}[tb]
  \centering
  \captionsetup{skip=3pt}% tighter image-to-caption gap, this float only
  \includegraphics[width=\linewidth]{figs/pf_qualitative_accv.png}
  \caption{Source images and each critic's best edit, with the LAION aesthetic score above each image, for images where AutoPolish helped most. The right column explains what AutoPolish changed and why it scores higher. AutoPolish makes the largest changes, such as adding a player behind a basketball, which no other critic proposes.}
  \label{fig:c4qual}
\end{figure}

\subsection{Ablations}
\label{sec:ablation}

\begin{table}[t]
  \caption{How each critic behaves when a safeguard of the loop is removed, where $^\ast$ marks the full loop. \textbf{(a)} Gain with and without the accept-if-better rule. \textbf{(b)} How far the final images move from the original with and without the similarity floor, where the last column gives the gain of the images that fall below a similarity of 0.78. Without the rule most of the gain disappears, and without the floor the critics that write their own instructions let their images drift away from the original.}
  \label{tab:abl-guard}
  \centering
  \fontsize{7.5}{9}\selectfont
  \setlength{\aboverulesep}{0pt}\setlength{\belowrulesep}{0pt}
  \renewcommand{\arraystretch}{1.1}
  \setlength{\tabcolsep}{2.4pt}
  \begin{tabular*}{\linewidth}{@{\extracolsep{\fill}}l|cc|c|cccc@{}}
    \toprule
    \multirow{3}{*}{Critic} & \multicolumn{2}{c|}{(a) Accept-if-better rule, gain $\uparrow$\rule{0pt}{2.4ex}} & \multicolumn{5}{c}{(b) Identity of the final images} \\
    \cline{2-8}
    & \multirow{2}{*}{with$^\ast$} & \multirow{2}{*}{without (95\% CI)} & w/ floor$^\ast$\rule{0pt}{2.3ex} & \multicolumn{4}{c}{w/o the floor} \\
    \cline{4-8}
    \rule[-0.8ex]{0pt}{0pt} & & & edit size\rule{0pt}{2.3ex} & edit size & below 0.78 $\downarrow$ & \shortstack{lowest\\similarity $\uparrow$} & gain \\
    \midrule
    static string\rule{0pt}{2.3ex} & 0.164 & \textbf{0.074} [0.040, 0.109] & 0.032 & 0.032 & \textbf{0\%} & \textbf{0.88} & n/a \\
    blind VLM & 0.155 & 0.008 [$-$0.032, 0.050] & 0.036 & 0.081 & 9\% & 0.03 & 0.588 \\
    \textbf{AutoPolish}\rule[-0.7ex]{0pt}{0pt} & \textbf{0.170} & 0.024 [$-$0.015, 0.061] & 0.049 & 0.108 & 12\% & 0.07 & 0.502 \\
    \midrule
    \textit{reward-only}\rule{0pt}{2.3ex}\rule[-0.7ex]{0pt}{0pt} & 0.187 & 0.062 [0.031, 0.097] & 0.040 & 0.040 & 0\% & 0.85 & n/a \\
    \bottomrule
  \end{tabular*}
\end{table}

\textbf{Accept-if-better rule.} The AutoPolish loop keeps a new edit only if the independent scorer, the LAION aesthetic predictor, rates it higher than the best image so far. We test the effect of this rule by removing it, so that the best candidate of each round replaces the current image even when it scores lower. Table~\ref{tab:abl-guard}a shows that without the rule the gain of AutoPolish falls from 0.170 to 0.024 and that of the blind VLM from 0.155 to 0.008, so that neither differs from zero any more, while the static string and the reward-only ceiling keep less than half of their gain. The rule is therefore essential, because the critics propose bold changes and the independent scorer decides which of them to keep.

\textbf{Identity preservation.} AutoPolish also rejects any candidate whose similarity to the original falls below 0.78, so that the edited image keeps its identity. We test what this similarity floor protects by removing it and measuring how far the final images move from the original. Table~\ref{tab:abl-guard}b shows that without the floor the edit size of our method more than doubles, from 0.049 to 0.108, and 12\% of its final images fall below a similarity of 0.78, with the lowest at 0.07. These 12 images gain 0.502 on average, about three times the gain of the full loop. The floor therefore keeps AutoPolish an editor rather than an image generator.

\begin{table}[t]
  \caption{How the gain of each critic changes with two settings of the loop, where $^\ast$ marks the full loop. \textbf{(a)} Gain with one to three candidates per round. \textbf{(b)} Gain at each similarity floor. Fewer candidates cost every critic a similar small amount, and the floor limits only the critics that write their own instructions.}
  \label{tab:abl-setting}
  \centering
  \fontsize{7.5}{9}\selectfont
  \setlength{\aboverulesep}{0pt}\setlength{\belowrulesep}{0pt}
  \renewcommand{\arraystretch}{1.1}
  \setlength{\tabcolsep}{2.4pt}
  \begin{tabular*}{\linewidth}{@{\extracolsep{\fill}}l|ccc|ccccc@{}}
    \toprule
    \multirow{2}{*}{Critic} & \multicolumn{3}{c|}{(a) Candidates per round, gain $\uparrow$\rule{0pt}{2.4ex}} & \multicolumn{5}{c}{(b) Similarity floor, gain $\uparrow$} \\
    \rule[-0.8ex]{0pt}{0pt} & 1 & 2 & 3$^\ast$ & none & 0.70 & 0.78$^\ast$ & 0.86 & 0.94 \\
    \midrule
    static string\rule{0pt}{2.3ex} & 0.130 & 0.150 & 0.164 & 0.164 & 0.164 & 0.164 & \textbf{0.164} & \textbf{0.147} \\
    blind VLM & 0.131 & 0.148 & 0.155 & 0.196 & 0.162 & 0.155 & 0.150 & 0.122 \\
    \textbf{AutoPolish}\rule[-0.7ex]{0pt}{0pt} & \textbf{0.141} & \textbf{0.159} & \textbf{0.170} & \textbf{0.215} & \textbf{0.176} & \textbf{0.170} & 0.159 & 0.121 \\
    \midrule
    \textit{reward-only}\rule{0pt}{2.3ex}\rule[-0.7ex]{0pt}{0pt} & 0.156 & 0.174 & 0.187 & 0.187 & 0.187 & 0.187 & 0.185 & 0.159 \\
    \bottomrule
  \end{tabular*}
\end{table}

\textbf{Candidates per round.} At each round the editor produces three candidates from different random seeds, and the loop keeps the one that the independent scorer, the LAION aesthetic predictor, rates highest. We test whether this selection, rather than the critic, produces the gain by keeping only the first one or two candidates of every round. Table~\ref{tab:abl-setting}a shows that with a single candidate AutoPolish still gains 0.141, which is 83\% of its gain with three, and stays ahead of the static string and the blind VLM at 0.130 and 0.131. Going from one to three candidates adds between 0.024 and 0.034 to every critic, so the selection among candidates explains only a small part of the gain.

\textbf{Similarity floor.} Since the floor protects the identity of the image, we also test how much the results depend on its value by varying it between 0.70 and 0.94 and by removing it. Table~\ref{tab:abl-setting}b shows that the floor only limits the critics that write their own instructions, since the static string and the reward-only ceiling keep almost the same gain at every floor up to 0.86, whereas the gain of AutoPolish rises to 0.215 without a floor, above the reward-only ceiling of 0.187. All of this extra gain comes from the images that lose their identity in Table~\ref{tab:abl-guard}b. Tightening the floor to 0.94 instead lowers every critic to between 0.121 and 0.159. The floor therefore trades gain against closeness to the original, and AutoPolish is the critic that presses hardest against it.

\section{Conclusion}

We presented AutoPolish, an automated image editor that replaces the single person guiding each round of editing with a synthetic audience, many copies of one vision-language model that each react as a different viewer, and keeps an edit only when an independent scorer rates it higher and the image still resembles the original. We showed that such an audience captures how real groups of viewers differ, and that in editing it improves images more than a generic critic on average while proposing three times as many distinct changes, with the independent scorer and the similarity floor turning these bolder proposals into gains that keep the image recognizable. Since the audience is only as faithful as its profiles and our gains are measured by an automatic scorer, the natural next step is to test with real viewers whether they prefer the edits their synthetic counterpart requests.

% TODO FINAL: add acknowledgements here (omitted for double-blind review).


\bibliographystyle{splncs04}
\bibliography{references}


%%%%%%%%%%%%%%%%%%%%%%%%%%%%%%%%%%%%%%%%%%%%%%%%%%%%%%%%%%%%
% ------------------------------ Appendix ------------------------------
%%%%%%%%%%%%%%%%%%%%%%%%%%%%%%%%%%%%%%%%%%%%%%%%%%%%%%%%%%%%
% Skipped when built through autopolish_main.tex (which defines \mainonly).
\ifdefined\mainonly\else
\newpage
\begin{center}
  \rule{\textwidth}{4pt}\\[0.16in]
  {\bfseries\fontsize{14}{17}\selectfont Appendix\\[5pt]}
  {\bfseries\fontsize{17}{21}\selectfont AutoPolish: Improving Visual Aesthetics\\ with Synthetic Audience Critique\\[0.14in]}
  \rule{\textwidth}{1pt}
\end{center}
\vspace{1.1em}

\appendix
\setcounter{section}{0}
\renewcommand{\thesection}{A\arabic{section}}
\setcounter{figure}{0}
\renewcommand{\thefigure}{A\arabic{figure}}
\setcounter{table}{0}
\renewcommand{\thetable}{A\arabic{table}}

% Keep floats (Table A1) from landing above the appendix title block.
\suppressfloats[t]

% Row struts for the appendix tables. Like the main-text tables, these drop
% booktabs' padding around the rules, so the first and last rows of each block
% carry a strut instead.
\newcommand{\tstrut}{\rule{0pt}{2.3ex}}
\newcommand{\bstrut}{\rule[-0.7ex]{0pt}{0pt}}

This appendix gives further detail behind the main paper. Appendix~\ref{sup:data} summarizes the four datasets, Appendix~\ref{sup:calib} explains why the audience's ratings need calibration and how well the calibration carries over between datasets, Appendix~\ref{sup:probe} describes the classifier that reads a rater's profile from the comments, and Appendix~\ref{sup:fig2} lists the exact values behind Figure~\ref{fig:persona}. Appendix~\ref{sup:robust} tests the audience against the most likely objections and Appendix~\ref{sup:fairness} measures whether it treats all kinds of viewers alike. Appendix~\ref{sup:c4qual} shows more AutoPolish edits, Appendix~\ref{sup:prompts} gives every prompt verbatim, and Appendix~\ref{sup:example} follows one image through one round of the loop. We close with future directions, ethics and reproducibility. Every confidence interval comes from 1{,}000 bootstrap resamples of images or raters, as stated with each result.

\section{Datasets}
\label{sup:data}
Table~\ref{tab:data} summarizes the four datasets described in the experimental setup. They differ most in what they record about their raters, which is richest on PARA and LAPIS and poorest on EVA, and this order matches where the synthetic audience works best in Section~\ref{sec:audience}.

\begin{table}[!h]
  \caption{The four datasets, with the number of images and ratings we use and what each records about its raters. PARA and LAPIS describe their raters in more detail than EVA, and the synthetic audience works best on those two.}
  \label{tab:data}
  \centering
  \fontsize{7.5}{9}\selectfont  % llncs aliases \footnotesize to \small
  \setlength{\aboverulesep}{0pt}\setlength{\belowrulesep}{0pt}
  \renewcommand{\arraystretch}{1.0}
  \begin{tabularx}{\linewidth}{l|l|cc|>{\raggedright\arraybackslash}X}
    \toprule
    Dataset\tstrut\bstrut & Content & Images & Ratings & Rater information \\
    \midrule
    PARA~\cite{yang2022para}\tstrut & photographs & 4{,}000 & 52k & age, gender, education, art and photography experience, Big-Five personality \\
    EVA~\cite{kang2020eva} & photographs from AVA & 4{,}070 & 137k & age, gender, region, photographic skill \\
    LAPIS~\cite{maerten2025lapis} & paintings & 4{,}000 & 90k & nationality, interest in art, age \\
    Rapidata~\cite{rapidata2024}\bstrut & AI-generated pairs & 898 pairs & 23k votes & country and language of each voter, across 120 countries \\
    \bottomrule
  \end{tabularx}
\end{table}

\section{Calibrating the audience's ratings}
\label{sup:calib}
We examine why the VLM's raw ratings need the calibration of Section~\ref{sec:method}, what it fixes, and whether a calibration fit on one dataset can be reused on another.

\begin{figure}[tbp]
  \centering
  \includegraphics[width=\linewidth]{figs/ax_calibration_accv.png}
  \caption{How the VLM uses the rating scale compared with people, and what calibration fixes. \textbf{(left)} Share of answers on the most common value and \textbf{(middle)} entropy of the ratings, for the VLM and people. \textbf{(right)} Error of a single rating and of a group average, before and after calibration, per dataset and on average. The VLM piles its answers onto one value, and calibration fixes the group average far more than one rating.}
  \label{fig:cal}
\end{figure}

Figure~\ref{fig:cal} (left and middle) shows that the VLM places 34 to 56\% of its answers on the most common value against 5 to 30\% for people, and that on the wider scales of EVA and LAPIS it also uses fewer values, with an entropy of 1.7 and 2.3 bits against 3.0 and 4.2. Because this habit is a fixed trait of the model rather than noise, averaging over a panel cannot remove it, but a calibration that maps each raw rating onto the human scale can.

Figure~\ref{fig:cal} (right) shows that calibration cuts the error of the group average by 38 to 66\% but that of a single rating by only 6 to 39\%, leaving a single rating two to three times further from the truth, as expected from Section~\ref{sec:ceiling}, where most of one person's rating is their own taste.

\begin{table}[tbp]
  \caption{Error of the audience's average rating when the calibration is fit on one dataset (columns) and applied to another (rows), next to the raw error. A fit on another dataset always beats the raw error but does less well than a fit on the dataset itself, shown in bold.}
  \label{tab:caltransfer}
  \centering
  \fontsize{7.5}{9}\selectfont
  \setlength{\aboverulesep}{0pt}\setlength{\belowrulesep}{0pt}
  \renewcommand{\arraystretch}{1.0}
  \newcolumntype{C}[1]{>{\hsize=#1\hsize\centering\arraybackslash}X}
  \begin{tabularx}{\linewidth}{l|C{1}|C{1}C{1}C{1}}
    \toprule
    \multirow{2}{*}{Evaluated on} & \multirow{2}{*}{Raw $\downarrow$} & \multicolumn{3}{c}{Calibration fit on $\downarrow$\rule{0pt}{2.4ex}\rule[-0.7ex]{0pt}{0pt}} \\
    \cline{3-5}
    \bstrut & & PARA\tstrut & EVA & LAPIS \\
    \midrule
    PARA\tstrut & 0.184 & \textbf{0.062} & 0.083 & 0.070 \\
    EVA & 0.103 & 0.094 & \textbf{0.063} & 0.085 \\
    LAPIS\bstrut & 0.145 & 0.086 & 0.100 & \textbf{0.072} \\
    \bottomrule
  \end{tabularx}
\end{table}

Table~\ref{tab:caltransfer} shows that a calibration fit on another dataset always beats the raw error, for example 0.070 against 0.184 on PARA with a fit on LAPIS, though it never matches a fit on the dataset itself. A fit on rated images can therefore be reused on unrated ones, such as the AI-generated images of Section~\ref{sec:generated}.

\section{Reading the persona from the comments}
\label{sup:probe}
Section~\ref{sec:persona} reports that the comments the VLM writes carry the persona it was given, even where its scores do not. Here we describe the classifier behind that result and the data it is trained on, together with two checks that rule out the most likely alternative explanations. Table~\ref{tab:probe} gives every number.

\begin{table}[!h]
  \caption{Recovering a rater's attribute from the VLM's comment alone, with and without a persona, under random folds and under folds split by rater, and the number of distinct comments per rater of an image. Art experience has four classes and the other two attributes three. The persona lifts the AUC well above chance under both kinds of fold, while the blind variant stays at chance.}
  \label{tab:probe}
  \centering
  \fontsize{7.5}{9}\selectfont
  \setlength{\aboverulesep}{0pt}\setlength{\belowrulesep}{0pt}
  \renewcommand{\arraystretch}{1.0}
  \newcolumntype{C}[1]{>{\hsize=#1\hsize\centering\arraybackslash}X}
  \begin{tabularx}{\linewidth}{ll|C{1}C{1}|C{1}C{1}|C{1}C{1}}
    \toprule
    \multirow{2}{*}{Dataset} & \multirow{2}{*}{Attribute} & \multicolumn{2}{c|}{Persona AUC $\uparrow$\rule{0pt}{2.4ex}\rule[-0.7ex]{0pt}{0pt}} & \multicolumn{2}{c|}{Blind AUC} & \multicolumn{2}{c}{Distinct comments $\uparrow$} \\
    \cline{3-8}
    \bstrut & & random\tstrut & by rater & random & by rater & persona & blind \\
    \midrule
    PARA\tstrut & art experience & \textbf{0.578} & \textbf{0.579} & 0.498 & 0.503 & \textbf{0.44} & 0.11 \\
    EVA & photographic skill & \textbf{0.701} & \textbf{0.704} & 0.511 & 0.513 & \textbf{0.24} & 0.08 \\
    LAPIS\bstrut & interest in art & \textbf{0.695} & \textbf{0.695} & 0.506 & 0.507 & \textbf{0.42} & 0.09 \\
    \bottomrule
  \end{tabularx}
\end{table}

\subsubsection{Classifier and data.} For each dataset, we pair every comment the VLM wrote while playing a rater with one attribute of that rater, namely art experience on PARA, photographic skill on EVA and interest in art on LAPIS, and we discard the score so that the classifier sees only the text. Numeric attributes are split into three groups of equal size, categorical ones keep their own levels, and any level with fewer than 500 comments is dropped. We draw 30{,}000 comments at random from each dataset (seed 0), out of 51.7k on PARA, 136.9k on EVA and 90.3k on LAPIS, which come from 436, 998 and 530 raters. Each comment is turned into TF-IDF weights over the 2{,}000 most frequent words, ignoring common English words and words that appear in fewer than five comments, and a multinomial logistic regression from scikit-learn with the default L2 penalty ($C{=}1$, at most 300 iterations) predicts the attribute from these weights. We score it with three-fold cross-validation and report the one-versus-rest ROC-AUC averaged over classes, where 0.5 is chance. We use AUC rather than accuracy because the classes are uneven, most of all on EVA, where always guessing the largest class would already be right two times in three while its AUC would stay at 0.5.

\subsubsection{Why a simple classifier.} We keep the classifier deliberately simple, since it sees only which words appear in a comment, knows nothing about the image and draws a straight boundary between classes. Whatever it finds is therefore carried by plain word choice rather than by a powerful model reading patterns into noise, so we treat its scores as a conservative estimate of how much of the persona the comments carry.

\subsubsection{The image is not the signal.} Raters with different attributes may have rated somewhat different images, and since the comments describe those images, a classifier could in principle guess the attribute from the image content alone. The blind variant tests this directly, because it sees the same images for the same raters and goes through the same pipeline but receives no persona. Its classifier stays at chance on every dataset, between 0.50 and 0.51, so the image content carries no usable trace of the attribute and the gain under the persona comes from the persona itself.

\subsubsection{The classifier does not recognize individual raters.} Random folds place comments from the same rater in both the training and the test data, and because all of a rater's comments come from the same persona card, the classifier could learn the phrasing of particular cards instead of the attribute. We therefore repeat the evaluation with folds split by rater, so that no rater appears on both sides, and refit the TF-IDF weights inside each fold. The scores do not move, at 0.579, 0.704 and 0.695 against 0.578, 0.701 and 0.695 with random folds, and the blind variant stays at chance. The classifier has therefore learned how a kind of viewer writes, and this carries over to raters it has never seen.

\subsubsection{Distinct comments.} For every image rated by at least ten people, we divide the number of distinct comments by the number of raters and average this ratio over images. With personas, the VLM writes three to five times as many distinct comments per image as the blind variant.

\section{Exact values behind Figure~\ref{fig:persona}}
\label{sup:fig2}
Table~\ref{tab:fig2} lists every number plotted in Figure~\ref{fig:persona}, with its 95\% confidence interval from 1{,}000 bootstrap resamples. For the left and right panels we resample images, since each value is an average over images, and for the middle panel we resample raters, because the comments written for one rater are not independent of each other. The AUC is computed within each of the three folds split by rater and then averaged, so each rater is scored only by a classifier that never saw them in training. The prior always predicts the overall average rating, and chance is the AUC of a classifier that guesses.

\begin{table}[tbp]
  \caption{Exact values behind Figure~\ref{fig:persona}, with 95\% bootstrap confidence intervals in brackets. The rating error and the distinct comments are computed over 4{,}000, 4{,}070 and 4{,}000 images on PARA, EVA and LAPIS, and the AUC over 30{,}000 comments per dataset from 436, 998 and 530 raters, where chance is 0.5. Calibration beats the prior on every dataset, and the persona raises both the AUC and the number of distinct comments on every dataset.}
  \label{tab:fig2}
  \centering
  \fontsize{7.5}{9}\selectfont
  \setlength{\aboverulesep}{0pt}\setlength{\belowrulesep}{0pt}
  \renewcommand{\arraystretch}{1.1}
  \newcolumntype{C}[1]{>{\hsize=#1\hsize\centering\arraybackslash}X}
  \begin{tabularx}{\linewidth}{>{\raggedright\arraybackslash}p{2.75cm}l|C{1}C{1}C{1}}
    \toprule
    Panel\tstrut\bstrut & Variant & PARA & EVA & LAPIS \\
    \midrule
    \multirow{3}{=}{Rating error $\downarrow$\newline(left panel)}\tstrut & raw & 0.184 [0.180, 0.188] & 0.103 [0.101, 0.105] & 0.145 [0.142, 0.147] \\
    & calibrated & \textbf{0.062} [0.060, 0.065] & \textbf{0.063} [0.062, 0.065] & \textbf{0.072} [0.070, 0.074] \\
    & prior\bstrut & 0.102 [0.099, 0.106] & 0.081 [0.079, 0.083] & 0.105 [0.103, 0.108] \\
    \midrule
    \multirow{2}{=}{Attribute AUC $\uparrow$\newline(middle panel)}\tstrut & persona & \textbf{0.579} [0.570, 0.587] & \textbf{0.704} [0.688, 0.728] & \textbf{0.695} [0.676, 0.713] \\
    & blind\bstrut & 0.503 [0.498, 0.509] & 0.513 [0.504, 0.530] & 0.507 [0.499, 0.515] \\
    \midrule
    \multirow{2}{=}{Distinct comments $\uparrow$\newline(right panel)}\tstrut & persona & \textbf{0.436} [0.429, 0.444] & \textbf{0.245} [0.241, 0.248] & \textbf{0.421} [0.416, 0.425] \\
    & blind\bstrut & 0.114 [0.111, 0.117] & 0.085 [0.083, 0.086] & 0.092 [0.090, 0.094] \\
    \bottomrule
  \end{tabularx}
\end{table}

\section{Robustness of the synthetic audience}
\label{sup:robust}
We check that the audience's group signal does not depend on how the model samples its answers, on the one rating we report or on what the image shows, and that the model does not recall the images or their ratings.

\begin{figure}[tbp]
  \centering
  \includegraphics[width=\linewidth]{figs/pf_breadth_category.png}
  \caption{Where the audience's group prediction holds. \textbf{(left)} Error of the calibrated audience against the prior for each of the 16 rated aspects of the three datasets, where a point below the diagonal means the audience wins. \textbf{(right)} Error of the calibrated audience by what the image shows, on PARA. The audience beats the prior on every aspect but one, and night scenes are the hardest content to judge.}
  \label{fig:breadth}
\end{figure}

\begin{table}[tbp]
  \caption{Results with greedy decoding ($t{=}0$) and sampling at $t{=}0.7$, where unanimous panels is the share of panels where every persona gives the same score, error drop is the fall in group error from one persona to twenty, and departure correlation is the agreement of group departures from Figure~\ref{fig:steer}. Sampling barely changes any result.}
  \label{tab:temp}
  \centering
  \fontsize{7.5}{9}\selectfont
  \setlength{\aboverulesep}{0pt}\setlength{\belowrulesep}{0pt}
  \renewcommand{\arraystretch}{1.0}
  \newcolumntype{C}[1]{>{\hsize=#1\hsize\centering\arraybackslash}X}
  \begin{tabularx}{\linewidth}{ll|C{1}C{1}|C{1}C{1}}
    \toprule
    \multirow{2}{*}{Dataset} & \multirow{2}{*}{Decoding} & Unanimous\rule{0pt}{2.4ex} & Error drop & Separation & Departure \\
    \rule[-0.8ex]{0pt}{0pt} & & panels & 1$\to$20 & $\uparrow$ & corr.\ $\uparrow$ \\
    \midrule
    PARA\tstrut & $t{=}0$ & 0.48 & 0.003 & $+0.044$ & $+0.37$ \\
    PARA & $t{=}0.7$ & 0.42 & 0.004 & $+0.040$ & $+0.39$ \\
    EVA & $t{=}0$ & 0.51 & 0.003 & $-0.085$ & $-0.24$ \\
    EVA & $t{=}0.7$ & 0.53 & 0.003 & $-0.080$ & $-0.27$ \\
    LAPIS & $t{=}0$ & 0.04 & 0.011 & $+0.166$ & $+0.39$ \\
    LAPIS\bstrut & $t{=}0.7$ & 0.03 & 0.006 & $+0.172$ & $+0.41$ \\
    \bottomrule
  \end{tabularx}
\end{table}

Table~\ref{tab:temp} shows that sampling at a temperature of 0.7 instead of always taking the most likely answer moves the separation between groups and the correlation of group departures by at most 0.03. It also shows that about half of the panels on PARA and EVA give the same score from every persona, which is why adding personas lowers the error so little there.

Figure~\ref{fig:breadth} shows that the calibrated audience beats the prior on 15 of the 16 rated aspects, not only on the overall score, and that its error depends on what the image shows, from 0.055 for animals to 0.089 for night scenes.

\begin{table}[tbp]
  \caption{Further checks against memory and bookkeeping, each on the dataset where it can be run. Every check comes out as expected for a model that judges the image rather than recalling its rating.}
  \label{tab:robust}
  \centering
  \fontsize{7.5}{9}\selectfont
  \setlength{\aboverulesep}{0pt}\setlength{\belowrulesep}{0pt}
  \renewcommand{\arraystretch}{1.0}
  \begin{tabularx}{\linewidth}{l|c|>{\raggedright\arraybackslash}X|c}
    \toprule
    Check\tstrut\bstrut & Dataset & What we measure & Value \\
    \midrule
    Artist fame\tstrut & LAPIS & rank correlation of an artist's number of ratings with the error & $+0.03$ \\
    Painting style & LAPIS & error gap between the easiest and the hardest style & $0.04$ \\
    Repeated ratings & LAPIS & group error under each of three ways of handling repeats & $0.072$ \\
    Rater agreement & LAPIS & correlation of rater disagreement with the group error & $-0.19$ \\
    Stated difficulty & EVA & correlation of the VLM's stated difficulty with rater disagreement & $+0.14$ \\
    Linked aspects\bstrut & EVA & agreement of the VLM's links between rated aspects with people's & $0.77$ \\
    \bottomrule
  \end{tabularx}
\end{table}

Table~\ref{tab:robust} shows no sign that the VLM recalls the images or their ratings, since artists rated more often are not predicted better and every other check is consistent with a model that judges the image itself. Its prompt also has no example images, so the overlap between EVA and AVA cannot leak ratings.

\section{Who the audience speaks for}
\label{sup:fairness}
We check whether the VLM treats all kinds of viewers alike, since a tilt toward viewers described one way would pass unseen into every average built on it. For each attribute of the persona cards we measure the calibration gap, the spread of the VLM's bias across the levels of that attribute after calibration, so that a gap of zero means every level is treated alike.

\begin{figure}[tbp]
  \centering
  \includegraphics[width=\linewidth]{figs/pf_bias.png}
  \caption{Calibration gap of the VLM for each attribute of the persona cards, which is the spread of its bias across the levels of that attribute, sorted from largest to smallest, with a dashed line at 0.05. The VLM is close to fair on most attributes but not on nationality (0.50), region (0.23) and gender on LAPIS (0.17).}
  \label{fig:bias}
\end{figure}

Figure~\ref{fig:bias} shows that the gap is at most 0.05 for personality, age, art and photography experience, and gender on PARA and EVA, but reaches 0.50 for nationality on LAPIS, the largest effect we measure in this paper, followed by 0.23 for region on EVA and 0.17 for gender on LAPIS. The VLM is therefore close to fair on the attributes that are usually audited but far less so on where the viewer comes from, and since aesthetic predictors also filter the images that generative models learn from, a bias of this kind could spread beyond evaluation. We therefore compare every between-group result in Section~\ref{sec:audience} with the blind variant, which removes such an offset, and make no claim that the audience tells countries apart, as our test between countries on generated images finds no effect.

\section{More qualitative examples}
\label{sup:c4qual}
Figure~\ref{fig:c4-qual-full} shows the five images where AutoPolish gains most over the static string, from which Figure~\ref{fig:c4qual} takes its second to fourth rows, and Figure~\ref{fig:c4-qual2} shows eight more images where AutoPolish reaches the highest score of all critics. As in the main paper, the right column explains each AutoPolish edit. In many rows AutoPolish changes what the image contains, for example by turning a bench red against a gray wall, replacing the ovens behind a chef with a bright kitchen, adding a truck to an empty snowy courtyard, or replacing the black sky behind a tower with a pale one, while the other critics mostly adjust light and color. In other rows its edit is as small as theirs and only deepens contrast or enriches color, which the explanations state plainly.

\begin{figure}[tbp]
  \centering
  \includegraphics[width=\linewidth]{figs/ax_qualitative_top5_accv.png}
  \caption{Source images and each critic's best edit, with the LAION aesthetic score above each image, for the five images where AutoPolish gains most over the static string. The right column explains what AutoPolish changed and why it scores higher. AutoPolish adds or changes objects, such as a red bench or a starry sky, while the other critics mostly adjust light and color.}
  \label{fig:c4-qual-full}
\end{figure}

\begin{figure}[p]
  \centering
  \includegraphics[width=\linewidth]{figs/ax_qualitative_more_accv.png}
  \caption{Eight more images where AutoPolish reaches the highest score, shown as in Figure~\ref{fig:c4-qual-full}. AutoPolish changes the setting or adds an object in four rows, such as the kitchen behind the chef and the truck in the courtyard, and makes smaller changes to light and color in the other four.}
  \label{fig:c4-qual2}
\end{figure}

% Defined before the paragraph that cites it, so it takes the float page
% right after Figure A5 instead of landing after the Prompts section starts.
\begin{figure}[p]
  \centering
  \includegraphics[width=\linewidth]{figs/ax_progression_accv.png}
  \caption{How AutoPolish improves nine images over its rounds, with the LAION aesthetic score above each image. The right column explains each change. In eight of the nine rows, half or more of the final gain is reached by the second round.}
  \label{fig:c4-prog}
\end{figure}

Figure~\ref{fig:c4-prog} follows nine images through the rounds of AutoPolish, from which Figure~\ref{fig:c4prog} takes three rows. Each row reads from left to right as the image improves, and the right column explains the change from the source to the last round, such as a Milky Way appearing over a night road or a plate of food becoming brighter. In eight of the nine rows, half or more of the final gain is reached by the second round, which matches the averages in Figure~\ref{fig:c4quant} (left), and the later rounds add smaller refinements. Some rows change only a little, and the explanations say so.

%%%%%%%%%%%%%%%%%%%%%%%%%%%%%%%%%%%%%%%%%%%%%%%%%%%%%%%%%%%%
% ---- Prompts and a worked example ----
%%%%%%%%%%%%%%%%%%%%%%%%%%%%%%%%%%%%%%%%%%%%%%%%%%%%%%%%%%%%

\section{Prompts}
\label{sup:prompts}
Every judgment in this paper comes from a frozen VLM given one of the prompt pairs below, and no weights are trained. The synthetic audience and the blind variant use the two framings shown here. Inside the editing loop, a further prompt turns the complaints the panel has gathered so far into a single instruction for the image editor.

\subsubsection{Synthetic audience.} The persona card is inserted verbatim at \pt{\{description\}}.
\begin{promptbox}{System prompt}
\ttfamily
You are role-playing as a person described as follows: \{description\}\\[2pt]
You are shown an image, as if it appeared in your social media feed. React the way this specific person actually would: let their personality, background, values, and taste -- not generic politeness -- decide whether they love it, are indifferent, or dislike it. Not everything deserves a high score; be honest and critical when the image doesn't match this person's taste. Do not mention that you are an AI or that you are role-playing.
\end{promptbox}
\begin{promptbox}{User prompt}
\ttfamily
As this exact person, look at this image. Give it an overall aesthetic score from 0 to 100 for how much YOU personally like it, then name the SINGLE change that would most improve it for your taste. Respond with ONLY a JSON object and nothing else, in exactly this form:\\[2pt]
\{"score": <integer 0-100>, "comment": "<one concrete improvement, imperative>"\}
\end{promptbox}

\subsubsection{Blind variant.} The same task with every reference to a persona removed.
\begin{promptbox}{System prompt}
\ttfamily
You are an impartial, general-purpose judge of photographic aesthetics.\\[2pt]
You are shown one photograph. Judge it on its own merits the way a typical viewer would, without adopting any particular person's perspective or taste. Not every photo deserves a good score; be honest when a photo is mediocre or disappointing. Do not mention that you are an AI.
\end{promptbox}
\begin{promptbox}{User prompt}
\ttfamily
Look at this image. Give it an overall aesthetic score from 0 to 100, then name the SINGLE most important change that would most improve it. Respond with ONLY a JSON object and nothing else, in exactly this form:\\[2pt]
\{"score": <integer 0-100>, "comment": "<one concrete improvement, imperative>"\}
\end{promptbox}

\subsubsection{From complaints to one instruction.} The current instruction is fed back in every round, so the complaints accumulate across rounds.
\begin{promptbox}{System prompt}
\ttfamily
You are a photo-editing assistant. You convert viewer feedback into ONE concrete, imperative image-edit instruction. Output only the instruction -- no preamble, no quotes, no explanation.
\end{promptbox}
\begin{promptbox}{User prompt}
\ttfamily
[The current edit instruction is: "<accumulated>".]\\[2pt]
Viewers looking at this image want these changes:\\
- <complaint 1>\\
- <complaint 2>\\
\ \ \dots\\[2pt]
Rewrite everything into ONE concrete imperative photo-edit instruction of at most 15 words that an image editor can apply. Output only the instruction.
\end{promptbox}

\subsubsection{Editor instructions shared by all critics.} The static string uses one fixed instruction in every round, while the other critics write their own. For every critic, the same request for a visible edit is appended to the instruction before it reaches the image editor.
\begin{promptbox}{Static string (fixed instruction)}
\ttfamily
Improve the overall aesthetic quality of this image.
\end{promptbox}
\begin{promptbox}{Visible-edit request (appended for every critic)}
\ttfamily
Make this a clearly visible edit, not a subtle one, while keeping the same subject and scene.
\end{promptbox}

\section{A worked example of one round}
\label{sup:example}
We follow one photograph of an open-plan office through one round of AutoPolish, starting from a LAION aesthetic score of 5.57 (Figure~\ref{fig:ex-edit}, left).

\subsubsection{The panel reacts.} The boxes below give the actual reactions of five of the ten personas on the panel, produced by the prompt of the synthetic audience in Appendix~\ref{sup:prompts}, where each persona returns a score out of 100 and the single change it would most want. The panel does not agree on one change, since two personas ask for warmer and more natural lighting, while the other eight ask for plants or artwork to make the desks more personal.

\begin{personabox}{Persona~1: male, 26 to 29, senior high school, art: competent \quad(score 60/100)}
``The lighting could be more natural to make the space feel more inviting.''
\end{personabox}
\begin{personabox}{Persona~5: male, 35 to 40, senior high school, art: expert \quad(score 60/100)}
``The lighting could be more natural to enhance the warmth and comfort of the space.''
\end{personabox}
\begin{personabox}{Persona~3: female, 30 to 34, university, art: proficient \quad(score 60/100)}
``Add some plants or artwork to the desks to add some life and personality to the space.''
\end{personabox}
\begin{personabox}{Persona~6: male, 35 to 40, junior college, art: beginner \quad(score 60/100)}
``Add some plants or artwork to the desks to make the space feel more welcoming and lively.''
\end{personabox}
\begin{personabox}{Persona~9: female, 18 to 21, junior college, art: beginner \quad(score 60/100)}
``Add some plants or artwork to the desks to make the space feel more welcoming and personal.''
\end{personabox}

\subsubsection{From complaints to an edit.} The loop gathers the ten complaints and turns them into one instruction of at most fifteen words with the prompt in Appendix~\ref{sup:prompts}, which merges the panel's two themes of lighting and personal touches. FLUX.1-Kontext applies this instruction to the original image, and the loop keeps the result because the independent scorer rates it higher while it stays close to the original. The kept edit adds warmer light, potted plants and framed artwork, which raises the score from 5.57 to 5.96 (Figure~\ref{fig:ex-edit}, right).

\begin{instrbox}{Complaints gathered from the ten personas}
$\bullet$~more natural, warmer lighting to make the space inviting \hfill (2 personas)\\
$\bullet$~add plants or artwork to the desks to add life and personality \hfill (8 personas)
\end{instrbox}
\begin{instrbox}{Instruction for the editor (at most 15 words)}
\itshape
Add natural lighting and plants/artwork to the desks to enhance warmth, comfort, and personalization.
\end{instrbox}

\begin{figure}[tbp]
  \centering
  \begin{subfigure}{0.40\linewidth}\centering
    \includegraphics[width=\linewidth]{figs/ex_office_src.jpg}\caption{source, score 5.57}\end{subfigure}\hspace{0.04\linewidth}
  \begin{subfigure}{0.40\linewidth}\centering
    \includegraphics[width=\linewidth]{figs/ex_office_edit.png}\caption{after one round, score 5.96}\end{subfigure}
  \caption{One photograph before and after one round of AutoPolish, with the LAION aesthetic score of each. The edit adds plants and framed artwork and warms the light, which are the changes the panel asked for, and raises the score from 5.57 to 5.96.}
  \label{fig:ex-edit}
\end{figure}

\section{Future directions}
\label{sup:future}
Our results come from a single frozen VLM, Qwen2-VL-7B-Instruct, so a first extension is to repeat the study with other models, across model families and across sizes within one family, to tell whether the synthetic audience is a property of the method or of this one model. Nothing in the method is specific to images either, so the same audience could react to other media such as audio, where datasets that record information about each listener would supply the persona cards and an audio editor would replace the image editor. Finally, the persona reaches the model only as text in the prompt, which limits how closely it can represent a viewer. More tailored ways of conditioning the model could close this gap, for example by learning part of the persona with low-rank adapters \cite{hu2022lora} attached to the frozen base model, with one adapter per attribute or group of viewers, since our results find the signal in groups rather than in single viewers.

\section{Ethics}
\label{sup:ethics}
Every dataset we use was collected and released for research by other groups, and we collected no data of our own. Our persona cards hold only attributes those datasets already publish, and a card describes a kind of viewer rather than a person, so we make no claim about any individual's taste. Biases in these records can carry into the judge, and the largest one we measure, a miscalibration by viewer nationality, is something anyone deploying such a system should check on their own population.

\section{Reproducibility}
\label{sup:repro}
We will publicly release the code, prompts and raw model outputs behind every number and figure in this paper. Nothing is trained, every model we use is frozen and publicly available, and the datasets can be obtained from their original authors. Appendix~\ref{sup:prompts} gives every prompt verbatim, and the experimental setup in the main paper describes the evaluation protocol.

\fi % end of the appendix (see \mainonly above)

\end{document}
